\documentclass[12pt,a4paper]{article}

% =====================================================
% PACOTES
% =====================================================
\usepackage[utf8]{inputenc}
\usepackage[T1]{fontenc}
\usepackage[brazil]{babel}
\shorthandoff{"}
\usepackage{geometry}
\usepackage{graphicx}
\usepackage{fancyhdr}
\usepackage{titlesec}
\usepackage[table]{xcolor}
\usepackage{hyperref}
\usepackage{setspace}
\usepackage{enumitem}
\usepackage{newtxtext,newtxmath}
\usepackage[most]{tcolorbox}
\usepackage{float}

\geometry{
    top=2.5cm,
    bottom=3cm,
    left=1.5cm,
    right=1.5cm
}

\onehalfspacing

% =====================================================
% CORES PERSONALIZADAS
% =====================================================
\definecolor{azulTitulo}{RGB}{30,60,90}
\definecolor{azulObjetivos}{RGB}{30,60,90}
\definecolor{cinzaCodigo}{RGB}{245,245,245}
\definecolor{vermelhoCodigo}{RGB}{180,50,50}
\definecolor{mostarda}{RGB}{204,153,0}

% =====================================================
% BOXES
% =====================================================
\newtcolorbox{boxAtividade}[1][]{
    breakable,
    enhanced,
    colback=azulTitulo!5,
    colframe=azulTitulo,
    boxrule=0pt,
    arc=6pt,
    outer arc=6pt,
    boxsep=0pt,
    left=14pt,
    right=14pt,
    top=16pt,
    bottom=14pt,
    title=#1,
    coltitle=white,
    colbacktitle=azulTitulo,
    fonttitle=\bfseries,
    attach boxed title to top left={
        yshift=-3mm,
        xshift=6mm
    },
    boxed title style={
        sharp corners,
        boxrule=0pt
    }
}

\newtcolorbox{boxCodigo}{
    colback=cinzaCodigo,
    colframe=vermelhoCodigo,
    boxrule=0.8pt,
    arc=2pt,
    left=6pt,
    right=6pt,
    top=6pt,
    bottom=6pt,
    borderline={0.5pt}{0pt}{vermelhoCodigo,dashed}
}

% =====================================================
% CABEÇALHO E RODAPÉ
% =====================================================
\pagestyle{fancy}
\fancyhf{}
\fancyhead[L]{
    % Tenta carregar a logo se o caminho estiver correto, senão ignora
    \IfFileExists{../../uc-topicos-avancados-latex/figures/garopaba_horizontal_marca2015_PNG.png}{
        \includegraphics[height=1.4cm]{../../uc-topicos-avancados-latex/figures/garopaba_horizontal_marca2015_PNG.png}
    }{IFSC - Garopaba}
}
\fancyhead[R]{
    \textbf{Manual do Professor}\\
    UC – Tópicos Avançados em Informática\\
    Técnico Integrado em Informática
}
\fancyfoot[L]{\small\textit{Manual de Conectividade - Roteiro 06}}
\fancyfoot[R]{\small Página \thepage}
\renewcommand{\footrulewidth}{2.4pt}
\setlength{\headheight}{45pt}

% =====================================================
% ESTILO DE SEÇÃO
% =====================================================
\titleformat{\section}
{\normalfont\Large\bfseries\color{azulTitulo}}
{\thesection}
{1em}
{}
[\vspace{0.3cm}\hrule]

\begin{document}

\section*{Manual do Professor: Conectividade e Monitoramento}

Este guia serve como referência rápida para configurar o ambiente do \textbf{Porto Central} antes do início das aulas, garantindo o acesso dos alunos e a supervisão do tráfego.

\begin{boxAtividade}[1. Pré-requisitos (Configuração NAT no VirtualBox)]
Para que o servidor Linux (VM) seja acessível pelo Mac (Host), certifique-se de que a VM está em modo \textbf{NAT} com as seguintes regras de \textbf{Redirecionamento de Portas} (Port Forwarding):

\vspace{0.3cm}
\centering
\begin{tabular}{|l|l|l|l|l|}
\hline
\rowcolor{azulTitulo!20} \textbf{Nome} & \textbf{Prot.} & \textbf{Porta Hosp.} & \textbf{Porta Conv.} \\ \hline
ssh & TCP & 2222 & 22 \\ \hline
vscode & TCP & 8080 & 8080 \\ \hline
\end{tabular}
\end{boxAtividade}

\section*{2. Passo a Passo de Inicialização}

Siga esta sequência exata no terminal do seu Mac:

\subsection*{Passo A: Estabelecer Conexão SSH e Túnel de Monitoramento}
Abra o terminal do Mac e conecte-se à VM puxando o painel do Ngrok para a sua máquina local:

\begin{boxCodigo}
\begin{verbatim}
ssh -L 4040:127.0.0.1:4040 ifsc@127.0.0.1 -p 2222
\end{verbatim}
\end{boxCodigo}

\subsection*{Passo B: Iniciar o Túnel Público (Ngrok)}
Dentro do terminal SSH (já logado na VM), inicie o túnel usando o seu domínio estático reservado:

\begin{boxCodigo}
\begin{verbatim}
ngrok http --domain=wobbly-subscript-onshore.ngrok-free.dev 8080
\end{verbatim}
\end{boxCodigo}

\section*{3. Como Supervisionar a Aula}

\begin{itemize}
    \item \textbf{Acesso dos Alunos:} Os alunos devem acessar a URL fixa via navegador:\\
    👉 \texttt{https://wobbly-subscript-onshore.ngrok-free.dev}
    
    \item \textbf{Monitoramento (Traffic Inspection):} Para ver em tempo real quais arquivos os alunos estão acessando, abra no navegador do seu Mac:\\
    👉 \texttt{http://localhost:4040}
\end{itemize}

\begin{tcolorbox}[colback=mostarda!10, colframe=mostarda, title=Dica de Ouro]
No painel do Ngrok (porta 4040), você consegue ver o histórico de todas as requisições HTTP. Se um aluno estiver com erro no código, você verá a requisição falhando (status 404 ou 500) ali na hora!
\end{tcolorbox}

\section*{4. Troubleshooting Rápido}

\begin{itemize}
    \item \textbf{SSH falhou:} Verifique se a VM está ligada e se a porta 2222 está configurada no VirtualBox.
    \item \textbf{Ngrok não inicia:} Verifique se o token de autenticação foi adicionado.
    \item \textbf{Página não carrega:} Verifique se o serviço do VS Code está rodando na VM (\texttt{sudo systemctl status code-server@ifsc}).
\end{itemize}

\end{document}
